% Part: lambda-calculus
% Chapter: church-rosser
% Section: beta-eta-reduction

\documentclass[../../../include/open-logic-section]{subfiles}

\begin{document}

\olfileid{lam}{cr}{be}

\olsection{$\beta\eta$-લઘુકરણ}

$\beta\eta$-લઘુકરણ ($\bered$) માટે ચર્ચ--રોસર
ગુણધર્મ લાગુ પડે છે. અહીં $\beredone$નો અર્થ
$\bredone$ અથવા $\eredone$માંથી કોઈ એકનું
એક-પગલાનું સંકોચન છે.
\sourcecorrection{OLLAM-043}{સ્થિર અંગ્રેજી મૂળ
$\beredone$ સંકેત વાપરે છે, પરંતુ આ પ્રકરણમાં
તેના સંબંધની સ્પષ્ટ વ્યાખ્યા આપતું નથી. અગાઉ
$\bered$ બંને પ્રકારના સંકોચનનું સ્વવાચક-પરંપરિત
સંવરણ છે; તેથી અહીં એક-પગલાના સંબંધને
$\bredone$ અને $\eredone$ના સંઘ તરીકે સ્પષ્ટ
કર્યો છે.}

\begin{lem}\ollabel{lem:one-par}
  $M \beredone M'$ હોય તો $M \beredpar M'$.
\end{lem}

\begin{proof}
  $M \beredone M'$ની !!{derivation} પર આગમનથી.
  મૂળ $\beta$-સંકોચન હોય તો
  \olref[pbe]{thm:refl} અને
  \olref[pbe]{defn:beredpar4} વાપરીએ; મૂળ
  $\eta$-સંકોચન, એટલે $M \eredone M'$, હોય તો
  \olref[syn][eta]{defn:beredone}ની તાજગી-શરત,
  \olref[pbe]{thm:refl} અને
  \olref[pbe]{defn:beredpar5} વાપરીએ.
  અમૂર્તીકરણ અથવા અનુપ્રયોગની અંદર સંકોચન
  થાય ત્યારે \olref[b]{lem:one-par}ની જેમ
  આગમનની ધારણા અને અનુક્રમે
  \olref[pbe]{defn:beredpar2} અથવા
  \olref[pbe]{defn:beredpar3} વાપરીએ.
  આથી દરેક એક-પગલાનું $\beta\eta$-સંકોચન
  સમાંતર પગલું છે.
  \sourcecorrection{OLLAM-044}{સ્થિર અંગ્રેજી મૂળ
  ``$M \bredone M'$ એ $\eta$-રૂપાંતરણથી'' એમ
  કહે છે; $\bredone$ તો ફક્ત $\beta$-સંકોચન
  દર્શાવે છે. અહીં $\eta$ માટે $\eredone$
  લીધો છે. મૂળની ટૂંકી દલીલમાં રચના-સુસંગત
  સંબંધના અંદરના કિસ્સા પણ ગાયબ હતા; તે
  આગમનથી ઉમેર્યા છે.}
\end{proof}

\begin{lem}\ollabel{lem:par-red}
  $M \beredpar M'$ હોય તો $M \bered M'$.
\end{lem}

\begin{proof}
  $M \beredpar M'$ની !!{derivation} પર આગમનથી.
  પહેલા ચાર નિયમો માટે
  \olref[b]{lem:par-red} જેવી દલીલ છે.

  છેલ્લો નિયમ \olref[pbe]{defn:beredpar5} હોય તો
  $M$ એ $\lambd[x][Nx]$ અને $M'$ એ $N'$, કોઈક
  $x$, $N$, $N'$ માટે, જ્યાં
  $x \notin FV(N)$ અને $N \beredpar N'$.
  પહેલાં $\eta$-રૂપાંતરણથી
  $\lambd[x][Nx]$નું $N$માં લઘુકરણ થાય.
  પછી આગમનની ધારણાથી $N \bered N'$ માટેની
  $\beredone$-પગલાંની શ્રેણી લાગુ કરીએ;
  તેથી $\lambd[x][Nx] \bered N'$.
\end{proof}

\begin{lem}\ollabel{lem:str}
  $\bered$ એ $\beredpar$ને સમાવતો સૌથી નાનો
  પરંપરિત સંબંધ છે.
\end{lem}

\begin{proof}
  \olref[b]{lem:str}ની દલીલ જેવી જ; શૂન્ય પગલાં
  માટે \olref[pbe]{thm:refl} વપરાય છે.
\end{proof}

\begin{thm}\ollabel{thm:cr}
  $\bered$ ચર્ચ--રોસર ગુણધર્મ ધરાવે છે.
\end{thm}

\begin{proof}
  \olref[dap]{thm:str}, \olref[pbe]{thm:cr} અને
  \olref{lem:str}થી દાવો મળે છે, જો
  \olref[pbe]{thm:cr}ની અગાઉ નોંધેલી સાબિતી-ખામી
  પૂરી કરવામાં આવે.
\end{proof}
\end{document}
